\documentclass[11pt,a4paper]{article}

% --- Encodage et langue ---
\usepackage[utf8]{inputenc}
\usepackage[T1]{fontenc}
\usepackage[french]{babel}
\usepackage{lmodern}

% --- Mise en page ---
\usepackage[margin=2.4cm]{geometry}
\usepackage{xcolor}
\usepackage{enumitem}
\usepackage{tcolorbox}
\usepackage[colorlinks=true,urlcolor=anstatvert,linkcolor=anstatvert]{hyperref}

% --- Couleurs ANStat (reprises du portail) ---
\definecolor{anstatvert}{HTML}{49655A}
\definecolor{anstatvertfonce}{HTML}{1F2D28}
\definecolor{anstatorange}{HTML}{F39323}
\definecolor{anstatgris}{HTML}{5C6C75}

\setlist[enumerate]{itemsep=3pt,topsep=4pt}
\setlist[itemize]{itemsep=3pt,topsep=4pt}

% --- En-tête du titre ---
\makeatletter
\renewcommand{\maketitle}{%
  \begin{center}
    {\small\color{anstatorange}\bfseries\MakeUppercase{StatLab \textperiodcentered{} ANStat}}\\[2pt]
    {\LARGE\bfseries\color{anstatvertfonce}\@title}\\[4pt]
    {\color{anstatgris}\@date}
  \end{center}
  \vspace{4pt}
  {\color{anstatorange}\hrule height 2pt}
  \vspace{10pt}
}
\makeatother

% --- Boîte « information » ---
\newtcolorbox{infobox}{
  colback=anstatorange!8, colframe=anstatorange,
  boxrule=1pt, arc=3pt, left=10pt, right=10pt, top=6pt, bottom=6pt
}

\newcommand{\service}[1]{\textbf{\color{anstatvertfonce}#1}}

\title{Se connecter à StatLab (ANStat)}
\date{Guide de l'utilisateur}

\begin{document}
\maketitle

\noindent
Ce guide explique comment accéder aux services de StatLab (environnements de travail,
stockage, tableaux de bord, requêtes SQL\dots) via le portail unique, avec votre
compte personnel.

\section*{Prérequis}
\begin{itemize}
  \item Être connecté au \textbf{VPN de l'ANSTAT}.
\end{itemize}

\section*{1.\quad Ouvrir le portail}
Dans votre navigateur, rendez-vous à l'adresse du portail StatLab fournie par
l'administrateur de la plateforme.
Vous êtes automatiquement redirigé vers la page de connexion.

\section*{2.\quad Se connecter (première connexion)}
\begin{enumerate}
  \item \textbf{Identifiant} : la partie située avant le \texttt{@} de votre adresse
        professionnelle. \\
        Exemple : pour \texttt{a.bakayoko@stat.plan.gouv.ci}, l'identifiant est
        \texttt{a.bakayoko}. \emph{(Votre adresse complète fonctionne aussi.)}
  \item \textbf{Mot de passe temporaire} :
        \begin{infobox}
          Il vous est \textbf{communiqué par l'administrateur de la plateforme} lors de
          la création de votre compte.
        \end{infobox}
  \item Le système vous demande aussitôt de \textbf{choisir votre propre mot de passe}.
        Faites-le : il est personnel et confidentiel.
  \item Vous arrivez sur le \textbf{portail StatLab}, avec les cartes de services.
\end{enumerate}

\section*{3.\quad Utiliser un service}
Cliquez sur une carte (par exemple \service{Environnements de travail},
\service{Stockage}, \service{Tableaux de bord}\dots). Le service s'ouvre et vous
\textbf{reconnaît automatiquement} : inutile de retaper votre mot de passe
(connexion unique).

\section*{Bon à savoir}
\begin{itemize}
  \item \service{Trino} (Requêtes SQL) : le navigateur affiche un avertissement de
        sécurité (certificat interne). C'est normal — cliquez sur « Avancé » puis
        « Continuer » pour accéder.
  \item \service{Environnements de travail} (JupyterHub) : au lancement, vous
        choisissez votre environnement (R, Python, VS~Code, RStudio\dots) et vous
        retrouvez vos fichiers habituels.
  \item \textbf{Un problème de connexion ?} Contactez l'administrateur de la plateforme.
\end{itemize}

\vspace{6pt}
{\color{anstatorange}\hrule height 1pt}
\vspace{4pt}
{\small\color{anstatgris}\itshape
Les connexions aux services sont protégées par un compte unique (authentification
centralisée). Conservez votre mot de passe personnel ; ne le partagez pas.}

\end{document}
